\chapter{Related Work and the State of the Art}
\label{chap:state_of_art_restructured}

This chapter reviews the literature that frames the present thesis: the integration of spatial regimes and temperature-derived descriptors into an adaptive AUV planner driven by uncertainty maps. The objective is not to catalogue all approaches to marine robotic sampling, but to identify the ideas that matter for a planner which starts from model-derived temperature and standard-deviation fields, converts them into operational reward maps, and evaluates whether additional structure extracted from the temperature field can improve sampling decisions.

The review is organised around six connected themes. Section~\ref{sec:sota_adaptive_sampling_auvs} introduces adaptive sampling with AUVs and the operational motivation for moving beyond static surveys. Section~\ref{sec:sota_model_driven_sampling} discusses model-driven ocean sampling, where numerical forecasts and ocean model products provide the prior information used by robotic systems. Section~\ref{sec:sota_uncertainty_planning} focuses on uncertainty-driven trajectory planning. Section~\ref{sec:sota_geostatistical_maps} reviews geostatistical uncertainty maps as a practical bridge between model products and AUV route selection. Section~\ref{sec:sota_regime_discovery} discusses spatial regime discovery and compact representations. Finally, Section~\ref{sec:sota_identified_gap} positions the thesis by identifying the opportunity to complement uncertainty maps with descriptors derived from spatial temperature regimes.

\section{Adaptive sampling with AUVs}
\label{sec:sota_adaptive_sampling_auvs}

Ocean observation is fundamentally constrained by the cost and sparsity of \emph{in situ} measurements. Ship-based surveys remain valuable, but they are expensive and episodic; fixed platforms provide persistence, but limited spatial coverage. Autonomous Underwater Vehicles (AUVs) address part of this limitation by collecting measurements along flexible trajectories, with a spatial resolution and deployment flexibility that are difficult to obtain from conventional observing assets. This motivation underpins the broader concept of autonomous ocean sampling networks, in which mobile platforms are integrated with ocean models, communication windows, and decision-making systems to collect observations where they are expected to be most informative \parencite{curtin2009aosn}.

Adaptive sampling differs from pre-planned surveying because the sampling plan is allowed to depend on the current state of knowledge. Rather than visiting a fixed grid, the vehicle is guided by a utility function derived from the estimated field, its uncertainty, or the presence of features of interest. Reviews of AUV adaptive sampling emphasise that this creates a sequential decision problem: the vehicle must balance information gain, travel cost, mission duration, safety, and the uncertainty associated with both the environment and the vehicle state \parencite{hwang2019auvreview}. This makes adaptive sampling a natural setting for informative path planning, where the route is selected not only to cover space, but to acquire measurements that improve a model or reduce uncertainty.

Several strands of work have contributed to this view. Early adaptive ocean sampling frameworks link sensing, modelling, data assimilation, and planning in a closed loop, where observations are used to update the representation of the environment and subsequent actions are chosen accordingly \parencite{lermusiaux2007adaptive}. In marine robotics, the same idea appears in feature-driven sampling, front tracking, plume monitoring, and targeted mapping, where vehicles are deployed to resolve structures that are spatially heterogeneous and temporally variable \parencite{das2012coordinated,fossum2018informationdriven,fossum2019phytoplankton}. Across these examples, the central lesson is that adaptive sampling requires a meaningful representation of value: the planner must know what makes one location more important than another.

For the present thesis, the relevant consequence is that an AUV planner can be seen as a mechanism that transforms environmental information into an operational decision. If the only available information is an uncertainty map, the planner will favour uncertain regions. If additional descriptors reveal thermal boundaries, gradients, or spatial regimes, then the planner can potentially distinguish between uncertainty that is spatially homogeneous and uncertainty that coincides with dynamically meaningful structure. This is the conceptual bridge from general adaptive sampling to the descriptor-enriched planner considered here.

\section{Model-driven ocean sampling}
\label{sec:sota_model_driven_sampling}

Model-driven ocean sampling uses numerical model products as an active component of the robotic decision process. Instead of treating models only as post-mission analysis tools, the model fields become inputs to the planner: they define the prior estimate of the environment, the uncertainty or error structure, and sometimes the operational constraints associated with currents, masks, or feasible domains. This idea is closely related to adaptive modelling and adaptive sampling, where observations and models are coupled in an iterative loop \parencite{lermusiaux2007adaptive}.

Recent work on ocean model-driven robotic exploration makes this connection explicit. In this perspective, ocean model forecasts and uncertainty products guide robotic sampling, while the observations collected by the robots are intended to improve subsequent model states or decision cycles \parencite{mendes2025oceanmodeldriven}. The value of this framing is operational: it avoids designing AUV missions in isolation from the ocean prediction system and instead treats robotic sampling as one component of a larger forecast--plan--sample--update pipeline.

\textcite{bernacchi2025closingloop} demonstrate this logic in a real-world AUV adaptive sampling setting, where model-derived uncertainty is converted into a planning objective and used to generate routes for data collection. The important contribution for the present thesis is the connection between an ocean model product, an uncertainty field, and a route planner. Once uncertainty is represented spatially, it can be masked, discretised, scored, and passed to a vehicle-routing formulation. This provides a practical template for using model outputs as planning inputs.

However, model-driven sampling also raises a question about what information should be transferred from the model to the planner. Model fields typically include physical variables such as temperature and salinity, as well as error or uncertainty proxies. A planner driven only by uncertainty may ignore spatial structure present in the physical field itself. Conversely, a planner driven only by physical gradients may over-sample strong features even where the model is already confident. The current thesis is positioned precisely at this interface: it asks whether uncertainty-driven planning can be enriched by descriptors extracted from the temperature field.

In the current implementation context, model products are obtained from CMEMS/IBI and interpolated to a high-resolution grid before deriving TEMP and STD maps. The existing bibliography does not yet contain a dedicated CMEMS/IBI reference, so a source should be added before final submission. \textbf{TODO: add and cite official CMEMS/IBI product documentation or an appropriate peer-reviewed reference.}

\section{Uncertainty-driven trajectory planning}
\label{sec:sota_uncertainty_planning}

Uncertainty is one of the most common signals used to drive adaptive sampling. If a model is uncertain in a region, then collecting observations there may reduce posterior uncertainty, improve reconstruction, or increase forecast skill. This idea appears in several forms: variance reduction, entropy reduction, mutual information, expected improvement, and reward maps derived from forecast error or model spread \parencite{hwang2019auvreview,lermusiaux2007adaptive}. Although the mathematical details differ, the planning logic is similar: candidate actions are scored by their expected contribution to information gain, and the route is constrained by vehicle feasibility.

The appeal of uncertainty-driven planning is that it provides a quantitative and interpretable objective. A spatial uncertainty map can be discretised into candidate sampling locations, and each location can be assigned a reward. The planner then selects a route that trades off reward against distance, time, energy, or other operational costs. Work on optimal synoptic surveys shows the value of moving beyond uniform grid sampling when uncertainty and decorrelation scales vary spatially \parencite{frolov2014grid}. In this sense, uncertainty maps provide a principled alternative to purely geometric coverage.

In robotic implementations, uncertainty-driven planning is often combined with myopic, receding-horizon, or vehicle-routing formulations. Myopic strategies are attractive because they are computationally light and responsive, but they may fail when travel time or future opportunities matter. Horizon-based and routing-based methods can reason over multiple targets, but require stronger assumptions about vehicle motion, candidate locations, and the value assigned to each location \parencite{hwang2019auvreview}. Flow-aware trajectory optimisation further shows that operational cost can depend strongly on time-varying currents, which affects feasibility and arrival times \parencite{aguiar2018trajectoryopt}.

The route-planning approach used by \textcite{bernacchi2025closingloop} is especially relevant because it transforms uncertainty into a routing reward in an operational AUV context. In that setting, a model-derived error or uncertainty field is filtered by mission constraints, discretised into candidate sites, and used by a prize-collecting routing formulation. The planner therefore does not need to reason directly over the full continuous model field; instead, it receives a structured planning problem built from model-derived rewards and constraints.

This thesis builds on that idea but changes the information content of the reward map. A baseline planner can be driven by STD alone, selecting trajectories that prioritise uncertain regions. An enriched planner can combine STD with descriptors computed from the temperature field, such as gradient magnitude or boundary scores. The central question is whether these descriptors help the planner choose trajectories that are not only uncertain, but also structurally informative with respect to spatial temperature regimes.

\section{Geostatistical uncertainty maps}
\label{sec:sota_geostatistical_maps}

Geostatistical methods provide an important practical route from sparse observations or model products to spatially explicit uncertainty maps. In contrast with highly complex physical simulations, geostatistical surrogates can produce predictions and associated uncertainty estimates at relatively low computational cost. This is attractive for AUV planning because the planner often requires a compact and frequently updated representation of where sampling is expected to be valuable.

\textcite{duarte2025geostatistical} address this need by constructing geostatistical uncertainty maps for efficient real-world AUV data collection. Their work is directly relevant because it treats uncertainty not merely as a diagnostic output, but as an operational planning input. The uncertainty field can be used to prioritise candidate sampling locations, while the map can be updated as measurements are collected. This gives a practical middle ground between full data assimilation and heuristic sampling: the planner receives an uncertainty product that is spatially structured, computationally manageable, and connected to the environmental field of interest.

From the perspective of the current thesis, geostatistical uncertainty maps are important for two reasons. First, they justify the use of STD or uncertainty-like fields as planner rewards. A spatial standard-deviation map is not just a visual diagnostic; it can be transformed into route utility. Second, they expose a limitation: uncertainty maps alone may not encode the full structure of the temperature field. A region with moderate uncertainty but a strong thermal boundary may be scientifically and operationally valuable, while a region with high uncertainty but weak structure may be less useful for understanding regime transitions.

This distinction motivates the combination of uncertainty maps with temperature-derived descriptors. The geostatistical literature supplies the planner-facing uncertainty component; regime and compact-representation literature supplies ideas for extracting structure from the physical field. The thesis sits between these two bodies of work by asking how descriptors derived from TEMP maps can modify or complement an STD-driven planner.

\section{Spatial regime discovery and compact models}
\label{sec:sota_regime_discovery}

Ocean fields are rarely homogeneous. Temperature, salinity, currents, and bathymetry can produce spatial regions with distinct behaviour, separated by fronts, gradients, or transitional boundaries. For adaptive sampling, these structures matter because measurements taken in one regime may not be equally informative for another. A planner that ignores regimes may sample uncertain areas but miss the boundaries that organise the field.

The literature on regime discovery and compact models provides tools for representing this structure. \textcite{fossum2020compact} propose compact models for adaptive sampling in marine robotics, using dictionary learning and sparse coding to represent high-dimensional observations in a lower-dimensional form. These compact representations can preserve relevant structure while reducing the dimensionality of the decision problem. The resulting codes can be clustered or otherwise analysed to identify recurring patterns, making them useful for regime-like structure discovery.

Related work on classification-driven adaptive sampling shows the value of probabilistic regime information. \textcite{ge2023watermasses} address adaptive AUV sampling for three-dimensional water-mass classification, where the vehicle collects measurements to improve classification over a spatial domain. The important point is not that this thesis adopts the same classification model, but that regime belief can become a planning signal. Ambiguous boundaries and uncertain class memberships are not merely errors; they can indicate where additional sampling is most informative.

Feature-driven sampling studies provide a similar lesson from another direction. Work on coordinated sampling of dynamic oceanographic features shows that vehicles can be directed towards fronts, plumes, or other evolving structures rather than treating the domain as a uniform field \parencite{das2012coordinated}. Information-driven robotic sampling further demonstrates how sampling can be shaped by the structure and variability of the target field \parencite{fossum2018informationdriven,fossum2019phytoplankton}. Together, these works support the idea that physical structure should influence the planning objective.

For the current thesis, the relevant regime information is extracted from gridded temperature maps rather than from labelled water-mass profiles or learned sparse codes alone. Descriptors such as \texttt{gradient\_magnitude} and \texttt{boundary\_score} provide lightweight ways to encode spatial structure in planner-ready form. They do not replace uncertainty maps; instead, they act as complementary layers that can increase the value of locations near thermal transitions or regime boundaries.

There is currently no dedicated reference in the project bibliography for the exact descriptors used in this implementation. If these descriptors are introduced as part of the thesis contribution, they can be described methodologically without an external citation. If they are based on a known image-processing or ocean-front-detection method, an appropriate reference should be added. \textbf{TODO: add references for temperature-front detection or image-gradient descriptors if the final method relies on established algorithms beyond standard finite differences.}

\section{Identified gap}
\label{sec:sota_identified_gap}

The reviewed literature establishes a clear trajectory. Adaptive sampling with AUVs has moved from static survey design towards decision-making based on environmental value \parencite{curtin2009aosn,hwang2019auvreview}. Model-driven ocean sampling connects this decision process to forecast products and operational ocean models \parencite{lermusiaux2007adaptive,mendes2025oceanmodeldriven}. Uncertainty-driven planning then turns model error, variance, or standard deviation into a spatial reward for route generation \parencite{frolov2014grid,bernacchi2025closingloop}. Geostatistical uncertainty maps make this approach practical by providing computationally efficient uncertainty products for AUV data collection \parencite{duarte2025geostatistical}. In parallel, regime-discovery and compact-model approaches show that ocean fields contain spatial structures and transitions that can be represented explicitly and used to inform sampling \parencite{fossum2020compact,ge2023watermasses}.

Despite these advances, there remains an opportunity at the interface between uncertainty-driven planning and spatial regime analysis. Existing planning pipelines commonly use uncertainty to decide where an AUV should sample, but uncertainty alone does not necessarily capture the structure of the underlying temperature field. A high-STD region may be valuable because the model is uncertain, but a thermal boundary or strong gradient may be valuable because it separates spatial regimes and can shape how observations should be interpreted. These are related but distinct sources of planning information.

The gap addressed by this thesis is therefore the limited integration of temperature-derived spatial regime descriptors into uncertainty-driven AUV trajectory planning. Rather than replacing STD-based planning, the thesis investigates whether descriptors such as \texttt{boundary\_score} and \texttt{gradient\_magnitude}, extracted from TEMP maps, can complement STD maps and produce an enriched planner that targets both uncertain and structurally informative regions. The resulting comparison between a baseline STD-driven planner and a descriptor-enriched planner provides a focused way to evaluate whether regime-aware information improves adaptive sampling in the Nazaré Canyon case study.

The Nazaré Canyon application is central to the thesis, but the current bibliography does not yet include a dedicated reference for the oceanographic relevance of the region. \textbf{TODO: add and cite one or more references on the Nazaré Canyon case study and its oceanographic dynamics.}
